% Chapter 5. Source pages PDF 115--148 (printed 107--140).
\section{Limits: definition and examples}
\begin{definition}[Binary product]
\label{def:lbc-product}
\source{leinster-basic-category-theory:page-0116}
\dcref{5.1.1}
A product of $X,Y$ is an object $P$ with maps $p:P\to X$, $q:P\to Y$ such
that for every $A$ and $f:A\to X$, $g:A\to Y$ there is a unique
$\langle f,g\rangle:A\to P$ with $p\langle f,g\rangle=f$ and
$q\langle f,g\rangle=g$.
\uses{def:lbc-category}
\end{definition}

\begin{remark}[Products are unique up to isomorphism]
\label{rem:lbc-product-uniqueness}
\source{leinster-basic-category-theory:page-0116}
\dcref{5.1.2}
Products need not exist, but when they do, the product cone is unique up to a
unique isomorphism. Strictly, the product is the object together with its two
projections; we often name only the object $X\times Y$.
\uses{def:lbc-product,def:lbc-isomorphism}
\end{remark}

\begin{definition}[Product of a family]
\label{def:lbc-family-product}
\source{leinster-basic-category-theory:page-0119}
\dcref{5.1.7}
For a set $I$ and objects $(X_i)$, a product is an object $P$ with projections
$p_i:P\to X_i$ such that maps $A\to P$ correspond uniquely to families
$(f_i:A\to X_i)_{i\in I}$.
\uses{def:lbc-product}
\end{definition}

\begin{definition}[Equalizer]
\label{def:lbc-equalizer}
\source{leinster-basic-category-theory:page-0120}
\dcref{5.1.11}
An equalizer of parallel maps $f,g:X\rightrightarrows Y$ is a map
$e:E\to X$ with $fe=ge$, universal among maps $h:A\to X$ satisfying $fh=gh$.
\uses{def:lbc-category}
\end{definition}

\begin{definition}[Pullback]
\label{def:lbc-pullback}
\source{leinster-basic-category-theory:page-0122}
\dcref{5.1.16}
A pullback of $f:X\to Z$ and $g:Y\to Z$ is an object $P$ with maps
$p:P\to X$, $q:P\to Y$, $fp=gq$, universal with this commutativity property.
\uses{def:lbc-product}
\end{definition}

\begin{definition}[Diagram and cone]
\label{def:lbc-diagram-cone}
\source{leinster-basic-category-theory:page-0126}
\dcref{5.1.18}
For a small category $I$, a diagram in $\mathcal A$ is a functor $D:I\to\mathcal A$.
A cone on $D$ with vertex $A$ is an objectwise compatible family
$\lambda_i:A\to D(i)$.
\uses{def:lbc-functor}
\end{definition}

\begin{definition}[Limit]
\label{def:lbc-limit}
\source{leinster-basic-category-theory:page-0126}
\dcref{5.1.19}
A limit of $D$ is a cone $(L,\lambda)$ such that for every cone $(A,\alpha)$
there is a unique map $u:A\to L$ with $\lambda_i u=\alpha_i$ for all $i$.
$\mathcal A$ has limits of shape $I$ when every such diagram has a limit and is
complete when it has all small limits.
\uses{def:lbc-diagram-cone}
\end{definition}

\begin{remark}[Limit cones and notation]
\label{rem:lbc-limit-notation}
\source{leinster-basic-category-theory:page-0127}
\dcref{5.1.20}
Maps into a limit object correspond bijectively to cones with that vertex. We
often call the vertex itself the limit and write it as $\limx D$, although the
whole limiting cone is the universal object.
\uses{def:lbc-limit,def:lbc-diagram-cone}
\end{remark}

\begin{definition}[Existence of limits]
\label{def:lbc-has-limits}
\source{leinster-basic-category-theory:page-0129}
\dcref{5.1.25}
A category has limits of shape $I$ if every diagram $I\to\mathcal A$ has a
limit, and has all limits if this holds for every small category $I$.
\uses{def:lbc-limit}
\end{definition}

\begin{proposition}[Products and equalizers suffice]
\label{prop:lbc-limits-products-equalizers}
\source{leinster-basic-category-theory:page-0129}
\dcref{5.1.26}
A category has all small limits if and only if it has all small products and all
equalizers.
\uses{def:lbc-family-product,def:lbc-equalizer,def:lbc-limit}
\notready
\end{proposition}
\begin{proof}
The source defers this proof to Exercise 5.1.38.
\end{proof}

\begin{definition}[Monomorphism]
\label{def:lbc-monic}
\source{leinster-basic-category-theory:page-0131}
\dcref{5.1.29}
A map $m:X\to Y$ is monic if $mf=mg$ implies $f=g$ for all parallel $f,g$.
\uses{def:lbc-category}
\end{definition}

\begin{lemma}[Monomorphism as equalizer]
\label{lem:lbc-monic-equalizer}
\source{leinster-basic-category-theory:page-0131}
\dcref{5.1.32}
A map $f:X\to Y$ is monic if and only if the square
\[
\begin{array}{ccc}
X&\xrightarrow{1_X}&X\\
\downarrow^{1_X}&&\downarrow^{f}\\
X&\xrightarrow{f}&Y
\end{array}
\]
is a pullback.
\uses{def:lbc-monic,def:lbc-pullback}
\end{lemma}
\begin{proof}
The source defers this proof to Exercise 5.1.41.
\end{proof}

\section{Colimits}
\begin{definition}[Colimit]
\label{def:lbc-colimit}
\source{leinster-basic-category-theory:page-0134}
\dcref{5.2.1}
A colimit of $D:I\to\mathcal A$ is a cocone $(C,\iota_i:D(i)\to C)$ universal
among cocones. A coproduct is a colimit of a discrete diagram.
\uses{def:lbc-limit,def:lbc-diagram-cone}
\end{definition}

\begin{definition}[Coproduct]
\label{def:lbc-coproduct}
\source{leinster-basic-category-theory:page-0135}
\dcref{5.2.2}
A coproduct (or sum) is a colimit of a discrete diagram.
\uses{def:lbc-colimit}
\end{definition}

\begin{definition}[Coequalizer]
\label{def:lbc-coequalizer}
\source{leinster-basic-category-theory:page-0136}
\dcref{5.2.7}
A coequalizer of $f,g:X\rightrightarrows Y$ is a map $q:Y\to Q$ with
$qf=qg$, universal among such maps.
\uses{def:lbc-colimit}
\end{definition}

\begin{remark}[Relations and generated equivalence relations]
\label{rem:lbc-generated-equivalence}
\source{leinster-basic-category-theory:page-0136}
\dcref{5.2.8}
A binary relation $R\subseteq A\times A$ is contained in a smallest equivalence
relation, obtained either by intersecting all equivalence relations containing
$R$ or by taking finite zigzags in the symmetric closure of $R$.
\uses{def:lbc-coequalizer}
\end{remark}

\begin{definition}[Pushout]
\label{def:lbc-pushout}
\source{leinster-basic-category-theory:page-0138}
\dcref{5.2.11}
A pushout of $X\xleftarrow f A\xrightarrow gY$ is a cocone
$X\to P\leftarrow Y$ universal among commutative cocones.
\uses{def:lbc-colimit}
\end{definition}

\begin{definition}[Epimorphism]
\label{def:lbc-epic}
\source{leinster-basic-category-theory:page-0141}
\dcref{5.2.17}
A map $e:X\to Y$ is epic if $fe=ge$ implies $f=g$ for all parallel maps out of
$Y$.
\uses{def:lbc-category}
\end{definition}

\section{Interactions with functors}
\begin{definition}[Preservation of limits]
\label{def:lbc-preserves-limits}
\source{leinster-basic-category-theory:page-0144}
\dcref{5.3.1}
A functor preserves limits of shape $I$ if it sends every limiting cone to a
limiting cone; it preserves all limits if this holds for every small $I$.
\uses{def:lbc-limit,def:lbc-functor}
\end{definition}

\begin{definition}[Creation of limits]
\label{def:lbc-creates-limits}
\source{leinster-basic-category-theory:page-0146}
\dcref{5.3.5}
A functor $F:\mathcal A\to\mathcal B$ creates limits of shape $I$ when every
limit cone for $FD$ lifts uniquely to a limit cone for $D$ whose image is the
given cone.
\uses{def:lbc-preserves-limits}
\end{definition}

\begin{lemma}[Creation implies preservation]
\label{lem:lbc-creates-preserves}
\source{leinster-basic-category-theory:page-0147}
\dcref{5.3.6}
Every functor that creates limits of shape $I$ preserves limits of that shape.
\uses{def:lbc-creates-limits,def:lbc-preserves-limits}
\notready
\end{lemma}
\begin{proof}
The source defers this proof to Exercise 5.3.12.
\end{proof}

\begin{remark}[Creation is stronger than preservation]
\label{rem:lbc-creation}
\source{leinster-basic-category-theory:page-0147}
\dcref{5.3.7}
Creation of limits includes a specified lift of each limiting cone, whereas
preservation only asserts that an already chosen limit remains limiting.
\uses{def:lbc-creates-limits,def:lbc-preserves-limits}
\end{remark}
